% GENERATED FILE. Edit canonical lessons, then run npm run generate:books.
% canonical-source-hash: fnv1a64:9c7cbab0004649cc
% canonical-lessons: KA-C171-yavagalu, KA-C171-agaga, KA-C171-dinanka, KA-C171-varantya, KA-C171-munce

\chapter{Always, Now and then, Date, Weekend, Before}
\label{ch:ka-always-now-and-then-date-weekend-before}
\hlchaptermodality{171}

\begin{chapteropening}
\textbf{By the end of this chapter:} \emph{I can say always, now and then, a date, the weekend and before.}

Complete the last lesson of chapter 171: I can say always, now and then, a date, the weekend and before.
\end{chapteropening}

\section[\texorpdfstring{\kn{ಯಾವಾಗಲೂ} (yāvāgalū)}{ಯಾವಾಗಲೂ (yāvāgalū)}]{\kn{ಯಾವಾಗಲೂ} (yāvāgalū) — always}

\label{lesson:KA-C171-yavagalu}



\begin{quote}
Before the new one: say the Kannada for a blister, then the Kannada for nutritious.
\end{quote}

\subsection*{You'll want to know: \kn{ಯಾವಾಗಲೂ}}
\textbf{\kn{ಯಾವಾಗಲೂ}} — \emph{yāvāgalū} — \textquotedblleft{}always\textquotedblright{}.

\begin{grammarlens}[title={What you've built}]
One more word for this chapter.
\end{grammarlens}

\subsection*{Your turn}
\begin{itemize}
\raggedright
  \item \emph{Say it:} \emph{yāvāgalū}
  \item \emph{Say it:} \emph{yāvāgalū}, once more
  \item \emph{Say it:} say \emph{yāvāgalū}
  \item \emph{Recall:} say the Kannada for a blister, then the Kannada for nutritious, then say \emph{yāvāgalū} again
\end{itemize}

\begin{tcolorbox}[breakable,colback=teal!4,colframe=teal!35!black,title={Before you move on}]
What does \kn{ಯಾವಾಗಲೂ} mean? (\textquotedblleft{}Always\textquotedblright{}.) Say it once more.
\end{tcolorbox}
\section[\texorpdfstring{\kn{ಆಗಾಗ} (āgāga)}{ಆಗಾಗ (āgāga)}]{\kn{ಆಗಾಗ} (āgāga) — now and then, from time to time}

\label{lesson:KA-C171-agaga}



\begin{quote}
Before the new one: say the Kannada for nutritious, then the Kannada for always.
\end{quote}

\subsection*{You'll want to know: \kn{ಆಗಾಗ}}
\textbf{\kn{ಆಗಾಗ}} — \emph{āgāga} — \textquotedblleft{}now and then, from time to time\textquotedblright{}.

\begin{grammarlens}[title={What you've built}]
One more word for this chapter.
\end{grammarlens}

\subsection*{Your turn}
\begin{itemize}
\raggedright
  \item \emph{Say it:} \emph{āgāga}
  \item \emph{Say it:} \emph{āgāga}, once more
  \item \emph{Say it:} say \emph{āgāga}
  \item \emph{Recall:} say the Kannada for nutritious, then the Kannada for always, then say \emph{āgāga} again
\end{itemize}

\begin{tcolorbox}[breakable,colback=teal!4,colframe=teal!35!black,title={Before you move on}]
What does \kn{ಆಗಾಗ} mean? (\textquotedblleft{}Now and then, from time to time\textquotedblright{}.) Say it once more.
\end{tcolorbox}
\section[\texorpdfstring{\kn{ದಿನಾಂಕ} (dināṅka)}{ದಿನಾಂಕ (dināṅka)}]{\kn{ದಿನಾಂಕ} (dināṅka) — a date (on the calendar)}

\label{lesson:KA-C171-dinanka}



\begin{quote}
Before the new one: say the Kannada for always, then the Kannada for now and then.
\end{quote}

\subsection*{You'll want to know: \kn{ದಿನಾಂಕ}}
\textbf{\kn{ದಿನಾಂಕ}} — \emph{dināṅka} — \textquotedblleft{}a date (on the calendar)\textquotedblright{}.

\begin{grammarlens}[title={What you've built}]
One more word for this chapter.
\end{grammarlens}

\subsection*{Your turn}
\begin{itemize}
\raggedright
  \item \emph{Say it:} \emph{dināṅka}
  \item \emph{Say it:} \emph{dināṅka}, once more
  \item \emph{Say it:} say \emph{dināṅka}
  \item \emph{Recall:} say the Kannada for always, then the Kannada for now and then, then say \emph{dināṅka} again
\end{itemize}

\begin{tcolorbox}[breakable,colback=teal!4,colframe=teal!35!black,title={Before you move on}]
What does \kn{ದಿನಾಂಕ} mean? (\textquotedblleft{}A date (on the calendar)\textquotedblright{}.) Say it once more.
\end{tcolorbox}
\section[\texorpdfstring{\kn{ವಾರಾಂತ್ಯ} (vārāntya)}{ವಾರಾಂತ್ಯ (vārāntya)}]{\kn{ವಾರಾಂತ್ಯ} (vārāntya) — the weekend}

\label{lesson:KA-C171-varantya}



\begin{quote}
Before the new one: say the Kannada for now and then, then the Kannada for a date.
\end{quote}

\subsection*{You'll want to know: \kn{ವಾರಾಂತ್ಯ}}
\textbf{\kn{ವಾರಾಂತ್ಯ}} — \emph{vārāntya} — \textquotedblleft{}the weekend\textquotedblright{}.

\begin{grammarlens}[title={What you've built}]
One more word for this chapter.
\end{grammarlens}

\subsection*{Your turn}
\begin{itemize}
\raggedright
  \item \emph{Say it:} \emph{vārāntya}
  \item \emph{Say it:} \emph{vārāntya}, once more
  \item \emph{Say it:} say \emph{vārāntya}
  \item \emph{Recall:} say the Kannada for now and then, then the Kannada for a date, then say \emph{vārāntya} again
\end{itemize}

\begin{tcolorbox}[breakable,colback=teal!4,colframe=teal!35!black,title={Before you move on}]
What does \kn{ವಾರಾಂತ್ಯ} mean? (\textquotedblleft{}The weekend\textquotedblright{}.) Say it once more.
\end{tcolorbox}
\section[\texorpdfstring{\kn{ಮುಂಚೆ} (muñce)}{ಮುಂಚೆ (muñce)}]{\kn{ಮುಂಚೆ} (muñce) — before, earlier}

\label{lesson:KA-C171-munce}



\begin{quote}
Before the new one: say the Kannada for a date, then the Kannada for the weekend.
\end{quote}

\subsection*{You'll want to know: \kn{ಮುಂಚೆ}}
\textbf{\kn{ಮುಂಚೆ}} — \emph{muñce} — \textquotedblleft{}before, earlier\textquotedblright{}.

\begin{grammarlens}[title={What you've built}]
One more word for this chapter.
\end{grammarlens}

\subsection*{Your turn}
\begin{itemize}
\raggedright
  \item \emph{Say it:} \emph{muñce}
  \item \emph{Say it:} \emph{muñce}, once more
  \item \emph{Say it:} say \emph{muñce}
  \item \emph{Recall:} say the Kannada for a date, then the Kannada for the weekend, then say \emph{muñce} again
\end{itemize}

\begin{tcolorbox}[breakable,colback=teal!4,colframe=teal!35!black,title={Before you move on}]
What does \kn{ಮುಂಚೆ} mean? (\textquotedblleft{}Before, earlier\textquotedblright{}.) Say it once more.
\end{tcolorbox}
